\documentclass[11pt]{scrartcl}
\usepackage[sexy]{evan}
\usepackage{answers}

\begin{document}
\title{Orders Solutions}
\author{Prakrit Gajurel}
\date{\today}
\maketitle

\section{Walkthrough Problems}

\begin{problem}
    Find all integers $n \ge 1$ such that $n$ divides $2^n - 1$.
\end{problem}

\begin{proof}
    We show that the answer is $n=1$ only. For the sake of contradiction, assume that some $n > 1$ works.
    Throughout this proof, $p$ represents an arbitrary prime number. Consider now any $p \mid n$.
    \begin{claim}
        $p \neq 2$
    \end{claim}
    \begin{proof}
        Obvious, as $p = 2 \implies 2 \mid n \mid 2^n - 1$ which is a contradiction.
    \end{proof}

    \begin{claim}
        If $p \mid n$, then $\text{ord}_p 2 \mid \gcd(n, p-1)$.
    \end{claim}
    \begin{proof}
        Since $p \mid n \mid 2^n - 1$, we get $2^n \equiv 1 \pmod{p}$ and so $\text{ord}_p 2 \mid n$. Moreover,
        Fermat gives $2^{p-1} \equiv 1 \pmod{p}$, so $\text{ord}_p 2 \mid p-1$ and the desired result follows.
    \end{proof}

    \begin{claim}
        For any $n \in \mathbb{Z}_{\ge 2}$, there exists a $p \mid n$ such that $\gcd(n, p-1) = 1$.
    \end{claim}
    \begin{proof}
        Choosing the smallest $p \mid n$ works(where $p$ can be $2$, for even $n$).
    \end{proof}

    \begin{claim}
        $\text{ord}_p 2 = 1$
    \end{claim}
    \begin{proof}
        Just choose the smallest $p \mid n$ which forces this.
    \end{proof}

    All of the above gives us that $2 \equiv 1 \pmod{p}$, which is a contradiction. Thus, $n=1$ is the only solution.
\end{proof}

\end{document}
